% Dies ist die Titelseite des Pflichtenhefts.
% Die in "<...>" sind zu ersetzen
% Die Ausgabe darf 1 Seite nicht �berschreiten, also ggf. Abst�nde anpassen
% Die Angabe in [...] gibt den Abstand nach der entsprechenden Zeile an. 


%----Stil dieser Seite--------------------------------------------------------------------
\thispagestyle{plain}			% Kopfzeile bleibt leer

%----Beginn der Titelseite----------------------------------------------------------------
\begin{titlepage}

%----zentrierte Ausrichtung �ber die gesamte Seite----------------------------------------		
\begin{center}

%----Titel des Praktikum (\praktikumTitel in newComments zu ver�ndern)--------------------
{\relsize{4}{\textbf{\textsc{\praktikumTitel}}}}\\[5ex]

%----Titel des Teilprojektes (\projektTitel in newComments ver�ndern)---------------------
{\relsize{3}{\textbf{\textsc{\projektTitel}}}}\\[5ex]

Softwareentwicklungspraktikum\\
Sommersemester 2008\\[6ex]

{\relsize{3}\so{\textbf{Pflichtenheft}}}\\[5ex]

%----eingebundenes Logo der TU------------------------------------------------------------
\includegraphics[scale=0.8]{bilder/carolo.jpg}\\[5ex]

%----Daten des Auftraggebers
Auftraggeber\\																
Technische Universit�t Braunschweig\\
<Name des Instituts>\\										
<Name des Institutsleiters>\\
<Stra�e und Hausnummer>\\
<Postleitzahl und Ort>\\[2ex]
Betreuer: <Name>\\[5ex]

% ----Tabelle der Praktikumsteilnehmer---------------------------------------------------
Auftragnehmer\\

\begin{tabular}{l<{\hspace{20mm}} l<{\hspace{30mm}}}\\	
	Name 									& 	E-Mail-Adresse\\			% Zeilen�berschift
		
	\hline										% Linie unterhalb der Zeilen�berschrift
	
	%----Nachfolgend alle Namen und E-Mail-Adressen der Teilnehmer einf�gen	
	<Name des Teilnehmers> &	<E-Mail-Adresse>\\
	<Name des Teilnehmers> &	<E-Mail-Adresse>\\
	<Name des Teilnehmers> &	<E-Mail-Adresse>\\
	...										 &	...	
	
\end{tabular}\\[2ex]

Braunschweig, <Abgabedatum>

\end{center}
\end{titlepage}